\begin{tabular}{lrrrrrrll}
\toprule
regime & encaminhados & taxa_encaminhamento & casos_detectados & casos_perdidos & recall & nns & regime_id & rotulo \\
\midrule
Regra SBD & 1421 & 0.840 & 115 & 4 & 0.966 & 12.357 & regra & Regra SBD \\
Random Forest · B_composto, recall equiparado (alvo 0.9768 no treino) & 1461 & 0.864 & 117 & 2 & 0.983 & 12.487 & equiparado & Random Forest · B_composto, recall equiparado (alvo 0,9768 no treino) \\
Random Forest · B_composto, ponto padrão (recall-alvo 0.80) & 685 & 0.405 & 96 & 23 & 0.807 & 7.135 & padrao & Random Forest · B_composto, ponto padrão (recall-alvo 0,80) \\
\bottomrule
\end{tabular}

\begin{flushleft}
\footnotesize Volume de encaminhamentos só é comparável entre estratégias que operam no mesmo recall. Quando os recalls realizados diferem, a comparação é por NNS, que normaliza por casos detectados — no regime de recall equiparado um modelo pode encaminhar mais pessoas que a regra em números absolutos e ainda assim ter NNS menor, por detectar mais casos.
\end{flushleft}
